\section{Consistent learning and description length: complete proofs}
\label{sec:annex-p15-occam-learning}

\begin{annexitem}{annex:p15.occam-finite}{thm:p15-occam-finite}{A union bound over every possible output}
\begin{proof}
Fix \(h\in H\) and let \(d_h=R_P(h,f)\). Agreement on
one observation has probability \(1-d_h\). Independence gives
probability \((1-d_h)^m\) of agreement on every observation.
For the fixed set \(H_{\rm bad}=\{h:d_h>\varepsilon\}\),
the event of a bad consistent hypothesis is a union of those
consistency events. If that set is nonempty, finite subadditivity gives
\[
 P^{\otimes m}(\text{a bad consistent hypothesis})
 \leq\sum_{h\in H_{\rm bad}}(1-d_h)^m
 <|H_{\rm bad}|(1-\varepsilon)^m
 \leq r(1-\varepsilon)^m.
\]
If the set is empty the probability is zero, still strictly
less than \(r(1-\varepsilon)^m>0\).
The bound covers every possible consistent selection
simultaneously; no independence of the learner's chosen
output from its sample is asserted or needed.

Let \(q=-\ln(1-\varepsilon)>0\). The condition
\(m\geq(\ln r+\ln(1/\delta))/q\) gives
\(r(1-\varepsilon)^m=r e^{-qm}\leq\delta\).
Taking an integer ceiling and requiring a nonempty sample
gives the displayed bound. For \(0<\varepsilon<1\),
\[
 q=\int_0^\varepsilon\frac{dt}{1-t}\geq\varepsilon,
\]
so replacing \(q\) by \(\varepsilon\) is sufficient.
The same inequality gives
\((1-\varepsilon)^m\leq e^{-m\varepsilon}\).
Each fixed-hypothesis agreement event is measurable because the
label functions are measurable and the label set is finite.
Finite unions preserve measurability. The weighted extension uses
only countable unions. The argument uses no particular \(P\) or target.
\end{proof}
\end{annexitem}

\begin{annexitem}{annex:p15.occam-threshold-example}{thm:p15-occam-threshold-example}{All errors, survivor events and sample sizes}
\begin{proof}
The target values at inputs \(0,1,2\) are \(0,0,1\).
Hypothesis \(h_1\) disagrees only at one, \(h_3\) only at two,
and \(h_0\) at zero and one. Their errors are \(1/3,1/3,2/3\),
respectively, so all three are bad at tolerance \(1/4\).
For \(h_1\) to survive, no sampled input may be one.
For \(h_3\) to survive, no sampled input may be two.
For \(h_0\) to survive, every sampled input must be two;
this is a subset of the first survivor event.
The union therefore consists exactly of the first two events.
Each has probability \((2/3)^m\), and their intersection
means that every input is zero, with probability \((1/3)^m\).
Inclusion-exclusion proves the formula.

For \(m\geq1\), its decrease from \(m\) to \(m+1\) is
\[
             \frac23\bigl((2/3)^m-(1/3)^m\bigr)>0.
\]
At nine and ten its values are
\((2^{10}-1)/3^9=1023/19683\) and
\((2^{11}-1)/3^{10}=2047/59049\).
The former is larger than \(1/20\) because \(20460>19683\);
the latter is smaller because \(40940<59049\).
Monotonicity proves minimality of ten.
For the generic bound, \(r/\delta=80\) and
\(q=\ln(4/3)\). The integer comparisons
\(4^{15}<80\cdot3^{15}\) and
\(4^{16}>80\cdot3^{16}\) prove its ceiling is sixteen.
Finally, input zero eliminates \(h_0\), input two eliminates
\(h_3\), and neither separates \(h_1\) from the target.
This proves the stated version space for that sample.
\end{proof}
\end{annexitem}

\begin{annexitem}{annex:p15.occam-code-count}{thm:p15-occam-code-count}{Counting all lengths and the prefix-tree inequality}
\begin{proof}
For an integer \(b\geq0\), there are \(2^j\) binary strings
of length exactly \(j\). Thus the number of strings of
length at most \(b\), including the empty one, is
\[
                \sum_{j=0}^b2^j=2^{b+1}-1.
\]
Every represented hypothesis has at least one such string,
and a string decodes to at most one hypothesis. Taking
\(b=\lfloor L\rfloor\) proves the arbitrary-code bound.

For a prefix-free family and an integer depth \(b\), each
code \(c\) of length at most \(b\) has \(2^{b-|c|}\)
extensions to binary words of length \(b\).
These descendant sets are disjoint: if one word extended
two codes, the shorter code would prefix the longer one.
There are only \(2^b\) words at depth \(b\), so
\[
 \sum_{\substack{c\in\mathcal C\\|c|\leq b}}2^{b-|c|}\leq2^b,
 \qquad
 \sum_{\substack{c\in\mathcal C\\|c|\leq b}}2^{-|c|}\leq1.
\]
Each summand is at least \(2^{-b}\), giving at most
\(2^b\) codes at lengths at most \(b\).
As \(b\) increases, the nonnegative partial sums increase
to the sum over every code; the limit is at most one.
Choosing one shortest code per hypothesis gives a subfamily
of distinct codes, preserving the inequality.
For arbitrary codes the three strings of lengths zero
or one show that \(2^L\) is not the general count
at \(L=1\). No prefix restriction was part of that count.
\end{proof}
\end{annexitem}

\begin{annexitem}{annex:p15.occam-compression}{thm:p15-occam-compression}{The entire compression-to-generalization implication}
\begin{proof}
Fix \(n,m\), the target of size at most \(n\), and \(P\).
Let \(L=n^c m^a\). Every possible output of the Occam
algorithm belongs to the finite set \(H_L\) of hypotheses
with code length at most \(L\). This set depends on
\(n,m\) and the fixed code system, not on the observed
sample. Because \(c\geq1,\ n,m\geq1\), \(L\geq n\),
so the target also belongs to \(H_L\).
The algorithm is consistent on every target sample.
Consequently failure is contained in the event of a
bad consistent member of \(H_L\).

For arbitrary codes, \(|H_L|<2^{L+1}\), so the preceding
finite union argument bounds failure by
\[
                2\exp(n^c m^a\ln2-qm).
\]
The stated second algebraic lower bound on \(m\) gives
\[
 m^{1-a}\geq\frac{2n^c\ln2}{q},
 \qquad n^c m^a\ln2\leq\frac{qm}{2}.
\]
The remaining bound gives
\(qm/2\geq\ln(2/\delta)\). Therefore failure is at most
\(2e^{-qm/2}\leq\delta\).
With prefix-free codes the leading factor two is absent,
and \(qm/2\geq\ln(1/\delta)\) suffices.
This proves both constant-level formulas, including
the case \(a=0\).

Since \(q\geq\varepsilon\), the sufficient sample size
is bounded by a constant multiple of
\[
 1+\frac{\ln(1/\delta)+\ln2}{\varepsilon}
       +\left(\frac{n^c}{\varepsilon}\right)^{1/(1-a)}.
\]
For \(n\geq1,\ 0<\varepsilon<1,\ 0\leq a<1\), the last
term is at least \(1/\varepsilon\), so it absorbs
the \(\ln2/\varepsilon\) term. This proves the stated
order uniformly for \(0<\delta<1\).
The fixed exponent \(1/(1-a)\) is finite.
Real fixed powers are bounded by integer powers on
arguments at least one, and \(\ln(1/\delta)\leq1/\delta\).
Thus the bound is polynomial in \(n,1/\varepsilon,1/\delta\).
The runtime requirement is already a premise of the
Occam algorithm, so combining it with this statistical
bound proves the source's polynomial learnability conclusion.
It does not give an algorithm for a class for which
the required consistent compression procedure is absent.
\end{proof}
\end{annexitem}

\begin{annexitem}{annex:p15.occam-minimum-approximation}{thm:p15-occam-minimum-approximation}{The sharper zero-exponent case}
\begin{proof}
When the output length is bounded by \(n^c\), apply the
finite union argument to the fixed class of codes of at
most that length. Arbitrary codes give failure at most
\[
              \exp((n^c+1)\ln2-qm).
\]
The stated lower bound on \(m\) makes this at most
\(\delta\). Prefix-free codes replace \(n^c+1\) by \(n^c\).
This direct use of the length bound avoids the factor-two
splitting used for general \(a\).

The actual target is a consistent candidate of length
at most \(n\), so \(b(S)\leq n\).
As \(n\geq1\), also \(b_+(S)\leq n\).
The assumed approximate minimizer therefore outputs
length at most \(n^c\). Its consistency and polynomial
runtime are premises, giving exactly the \(a=0\)
Occam condition. An exact minimum-length selector
outputs length \(b(S)\leq b_+(S)\), the case \(c=1\),
provided it satisfies the separate runtime condition.
\end{proof}
\end{annexitem}

\begin{annexitem}{annex:p15.occam-weighted}{thm:p15-occam-weighted}{A countable union with a fixed allocation}
\begin{proof}
For each \(h\), set
\[
                   t_h=\frac{\ln(1/w_h)+\ln(1/\delta)}m.
\]
The weights sum to at most one and are positive, so
\(w_h\leq1\) and \(t_h>0\).
If \(t_h\geq1\), error larger than \(t_h\) is impossible.
If \(0<t_h<1\) and \(R_P(h,f)>t_h\), the probability
that this fixed \(h\) is consistent is less than
\[
             (1-t_h)^m\leq e^{-mt_h}=\delta w_h.
\]
If its error is already at most \(t_h\), that hypothesis
cannot violate the conclusion. Countable subadditivity
therefore bounds the probability that any hypothesis
is both consistent and in violation by
\(\sum_h\delta w_h\leq\delta\).

For a prefix-free code, the preceding Kraft proof
justifies \(w_h=2^{-\ell(h)}\).
For arbitrary codes, choose one shortest representative
for each \(h\); these representatives are distinct.
At length \(j\) there are at most \(2^j\) of them. Hence
\[
       \sum_h2^{-2\ell(h)-1}
           \leq\sum_{j=0}^\infty2^j2^{-2j-1}
           =\sum_{j=0}^\infty2^{-j-1}=1.
\]
Substituting the two weight formulas proves both
description-length bounds. The proof fixes these
weights before the sample and performs no posterior
normalization over the hypotheses.
\end{proof}
\end{annexitem}

\begin{annexitem}{annex:p15.occam-boundaries}{thm:p15-occam-boundaries}{Three explicit changes of experiment}
\begin{proof}
For dependence, use \(X=\{0,1\}\), \(P(0)=P(1)=1/2\),
target \(f=0\), and \(h(x)=x\). Draw one \(Z\sim P\)
and set \(X_i=Z\) for every \(i\). Each marginal remains
\(P\), but consistency of \(h\) occurs exactly when
\(Z=0\), with probability \(1/2\). Independent inputs
would give \(P(X_1=\cdots=X_m=0)=2^{-m}\).
The fixed class \(\{f,h\}\) contains the target in both
experiments; the independence assumption is what changed.

For distribution change, use the same functions, but
train with input law \(\delta_0\) and assess future
error under \(\delta_1\). Every training observation
has correct label zero for \(h\), but its future error
is one. The independent-sample theorem still gives
the correct error zero under the training law; it
never claimed a guarantee under the different law.

For the adaptive class, let the target be identically
zero on an \(N\)-point domain. Let \(O(S)\) be the
set of observed inputs. Then \(|O(S)|\leq m\) and
\(h_S(x)=\mathbf1_{\{x\notin O(S)\}}\) is consistent.
Under uniform input probability its error is
\[
                  \frac{N-|O(S)|}{N}\geq1-\frac mN.
\]
The class containing the target and this chosen
hypothesis has at most two elements, but it varies
with \(S\). The earlier finite union must instead
range over a class fixed in advance that contains
every possible output. In the unrestricted binary
class this means counting up to \(2^N\) functions,
not merely the two named after seeing the data.
If two records of the same input have different
labels, no function on \(X\) agrees with both,
which also proves the stated empty-version-space
boundary under erroneous labels.
\end{proof}
\end{annexitem}

\begin{annexitem}{annex:p15.occam-unseen-labels}{thm:p15-occam-unseen-labels}{Averaging over all binary targets}
\begin{proof}
For this proof only, choose a random target \(F\) uniformly
from the \(2^N\) binary functions, independently of the
sampled inputs and the learner's internal randomization.
Equivalently, its \(N\) labels are independent fair bits.
Condition on the sampled inputs, their observed labels,
and any internal randomization. At an unobserved point,
the target label is still a fair bit, independent of
the learner's predicted binary value. Its conditional
error probability is therefore \(1/2\).

There are at least \(N-m\) unobserved points when
\(m\leq N\); errors at observed points are nonnegative.
Conditional expected uniform-input error is thus at
least \((N-m)/(2N)\). Taking expectation preserves
that bound. If \(m>N\), the displayed bound is
nonpositive and follows from nonnegativity.

Finally, the unconditional expected error is the
average, over the finite set of fixed targets, of
their expected errors under data and learner randomness.
At least one target must have expected error at
least this average lower bound. The random choice
of \(F\) was used to prove existence of that fixed
target, not inserted into the hypotheses of the
distribution-free upper bound.
\end{proof}
\end{annexitem}
